% AMS 560 midterm report. Build: latexmk -pdf midterm-report.tex (page limit: 2).
\documentclass[11pt]{article}

\usepackage[T1]{fontenc}
\usepackage{mathptmx}
\usepackage[letterpaper,margin=0.8in]{geometry}
\usepackage{booktabs}
\usepackage{tabularx}
\usepackage{array}

\setlength{\parindent}{0pt}
\setlength{\parskip}{4pt}
\renewcommand{\arraystretch}{1.12}
\newcolumntype{L}{>{\raggedright\arraybackslash}X}

% Compact section headings without extra packages.
\makeatletter
\renewcommand{\section}{\@startsection{section}{1}{0pt}%
  {6pt plus 2pt minus 1pt}{2pt}{\normalfont\large\bfseries}}
\makeatother

\pagestyle{plain}

\begin{document}

\begin{center}
  {\Large\bfseries Path-Aware Governance for High-Volume AI-Agent Workflows}\\[3pt]
  {\normalsize AMS 560 Project \textemdash{} Midterm Report}\\[1pt]
  % TODO(team): add team member names.
  {\small Team members: \rule{0pt}{0pt}}
\end{center}

\section{Objective and Project Description}

The project asks whether \emph{path-aware} governance, which judges each proposed agent action
in light of the earlier actions in the same workflow and across cooperating agents, catches
policy violations that per-action rules miss, and what it costs in latency, throughput and
state consistency. We compare three engines on one shared event format: \textbf{no governance}
(allow everything), \textbf{stateless} (rules on single actions: dangerous actions, a
per-transfer limit, unauthorized actors), and \textbf{path-aware} (stateless rules plus
multi-step rules for cross-agent data exfiltration, privilege laundering through delegation,
and split-transfer evasion). A Kafka pipeline runs the engines with several workers under
strong (workflow-partitioned) and weak (agent-partitioned, replicated) state.

The first phase evaluated the engines on a seeded, template-generated corpus of 150 labeled
workflows. Because those workflows and the policies share an author, we treat that result as
an engine sanity check, not as evidence about real agents. The second phase, now under way,
puts a large language model (LLM) agent in a simulated environment where every tool call
passes through the governance engine, so that detection and cost can be measured on behavior
we did not script.

\section{Timeline and Expected Deliverables}

\begin{tabularx}{\textwidth}{@{}l L l@{}}
\toprule
\textbf{Phase} & \textbf{Deliverable} & \textbf{Status} \\
\midrule
1 & Event schema, policy definitions, three governance engines & Complete \\
2 & Seeded synthetic workflow corpus with blind-labeling sample & Labels pending review \\
3 & Kafka pipeline (actions, decisions, state topics; replay; partitioning) & Complete \\
4 & Strong and weak multi-worker shared state & Complete \\
5 & Benchmarks (accuracy, stale reads, latency, throughput) and dashboard & Complete \\
6 & Agent test harness: simulated tools, governance gateway, provenance tracking & Complete \\
7 & LLM agent loop and single-model pilot (GPT-5.6 Luna, medium reasoning) & In progress \\
8 & Agent scenario suite, outcome oracle, repeated-run experiment design & Started \\
9 & Final validation, reproducibility check, final report & Planned \\
\bottomrule
\end{tabularx}

\section{Progress on Each Deliverable}

{\small
\begin{tabularx}{\textwidth}{@{}p{0.27\textwidth} L@{}}
\toprule
\textbf{Deliverable} & \textbf{What is done and what remains} \\
\midrule
Engines (1) & All three engines implemented behind one interface. The path-aware engine now
records every attempted action with a status (blocked, allowed, executed, failed); cumulative
limits count only executed transfers, so a blocked attempt no longer inflates later totals. \\
Workflow corpus (2) & 150 reproducible safe, violating and near-miss workflows (seed 560).
Remaining: independent two-reviewer labeling of the prepared 20-workflow sample. \\
Kafka pipeline and state (3, 4) & Complete. With four workers, workflow partitioning gave
0\% stale reads and 100\% recall; agent partitioning with replicated state gave about 25\%
median stale reads and 57--59\% median recall. \\
Benchmarks and dashboard (5) & Complete. On the synthetic corpus, stateless recall is 23.3\%
and path-aware recall is 100\% with no false positives (by construction, see Section~1).
In-process path-aware throughput was about 340{,}000 events/s at concurrency 10 and 207{,}000
at concurrency 100. \\
Agent harness (6) & Complete. A simulated world (data sources, files, accounts, permissions)
and eight tools; a gateway that builds each event from world state rather than from the agent's
claims and blocks or runs the call; lineage-based provenance so data-flow rules work without
agent cooperation. Eleven scripted scenarios run under all three levels with an outcome check
based on what actually happened in the world. Building it exposed two engine gaps that we
fixed (a harmless read before a PII read hid exfiltration; blocked attempts counted toward
transfer totals) and three that we document as known limits (merging two PII sources;
one agent splitting a payment; a legitimate over-limit payment with no approval path). \\
LLM pilot (7) & A model-neutral agent loop is built and tested with a scripted model. A
connection through the Codex CLI with our tools served over MCP was tested live: governed tool
calls work and built-in shell and web tools can be disabled, but Codex sub-agents each start a
separate copy of the tool server, which splits the workflow history. \\
Scenario suite and design (8) & Groundwork only: outcome-based harm check, basic metrics
(blocks, retries after a block, task completion), and offline replay of ungoverned runs through
the other engines. \\
\bottomrule
\end{tabularx}}

\section{Plan for Unfinished Deliverables}

\textbf{Corpus labels (2).} Two members label the 20-workflow blind sample using only the
written policies; we compare against each other and the generated labels and revise ambiguous
policies or templates before freezing them.

\textbf{LLM pilot (7).} Choose the connection: either fix the Codex route with one shared tool
server per run, or call the model through the OpenAI API with the key read from an environment
variable (never stored in the repository). Then run the three draft pilot tasks with one model
to debug the loop and the outcome check before any full experiment.

\textbf{Scenario suite and design (8).} Write scenarios where a violation is possible but never
instructed (benign, tempting, adversarial, multi-agent handoff), with some written by a member
who did not write the rules and a held-out set never used to tune rules. Freeze policies and
pre-register hypotheses and metrics, then run repeated trials per scenario and governance level
and report confidence intervals, task completion and benign-task blocking alongside recall.

\textbf{Final validation (9).} A second member reproduces the full setup on another machine
(dependencies, Kafka, tests, benchmarks, harness). The final report will state results on
synthetic workflows and on agent runs separately and list the known limits above.

\end{document}
